\section{\QSCver{legacy.37.0} Migration Map and Backtranslation Contract for the shifted QAM Shift}
\label{sec:v3370-migration-map-backtranslation-contract}

\begin{quote}\small
\noindent\textbf{Status.}
This section is a migration-map layer only. It prepares the later shifted global
QAM Hilbert-shift, but does not perform that global migration. The local
HC-boundary activation shift introduced in Version~3.36.0 remains the first
structural shift:
\[
 0_{\mathrm{mask}}\longrightarrow 1_{\mathrm{mask}}.
\]
The present layer records how a later shifted line may move the whole QAM tag
language without losing the historical release line Active Core, its router/readout separation,
or the Historical Witness.
\end{quote}

\subsection{Separation of the two shifts}
\label{subsec:v3370-two-shift-separation}

\begin{definition}[Local boundary shift versus global language shift]
\label{def:v3370-local-boundary-shift-versus-global-language-shift}
The Companion distinguishes two shifts.
\begin{enumerate}[leftmargin=*]
 \item \emph{Shift A}, the local HC-boundary activation shift, is the
 already-admitted historical release line structural reading
 \[
 0_{\mathrm{mask}}\longrightarrow 1_{\mathrm{mask}}.
 \]
 It says that the transition from no active readout to the first active
 readout requires the Heisenberg Compensator as a minimal boundary/readout
 normalizer.

 \item \emph{Shift B}, the global shifted QAM Hilbert-shift, is the future
 language migration
 \[
 k\longmapsto k+1
 \]
 applied uniformly to inhabited QAM tags so that tag $0$ becomes globally
 reserved for the empty or no-active-kind marker.
\end{enumerate}
\end{definition}

\begin{remark}[Why Shift A belongs to the legacy preparation line and Shift B belongs to shifted]
\label{rem:v3370-shift-a-legacy-shift-b-shifted}
Shift A is structural and local: it explains why HC appears internally at the
first active readout boundary. Shift B is linguistic and global: it changes
the whole QAM namespace. Therefore Shift A belongs to the legacy preparation
line, while Shift B must be treated as the first major shifted migration.
\end{remark}

\subsection{Old and new QAM namespaces}
\label{subsec:v3370-old-new-qam-namespaces}

\begin{definition}[Pre-shift and shifted tag namespaces]
\label{def:v3370-pre-shift-shifted-tag-namespaces}
Let $\mathcal N_{3}$ denote the inherited pre-shift QAM tag namespace used through
Version~3.36.0. Let $\mathcal N_{4}$ denote the future shifted namespace. The
planned migration map is a partial translation
\[
 \Sigma_{4}\colon \mathcal N_{3}^{\mathrm{inh}}\longrightarrow
 \mathcal N_{4}^{\mathrm{inh}},
 \qquad
 \Sigma_{4}(k)=k+1,
\]
defined on inhabited tags. The target tag $0\in\mathcal N_{4}$ is reserved for
the empty or no-active-kind marker and is not the image of an inhabited
pre-shift tag.
\end{definition}

\begin{definition}[Backtranslation map]
\label{def:v3370-backtranslation-map}
The formal backtranslation associated with $\Sigma_{4}$ is the partial map
\[
 \mathrm{B}_{4}\colon \mathcal N_{4}^{\mathrm{inh}}\longrightarrow
 \mathcal N_{3}^{\mathrm{inh}},
 \qquad
 \mathrm{B}_{4}(a)=a-1
\]
for inhabited shifted tags $a\geq 1$. The value $0$ has no inhabited
backtranslation; it is read only as the empty/no-active-kind marker.
\end{definition}

\begin{lemma}[Round-trip on inhabited tags]
\label{lem:v3370-round-trip-inhabited-tags}
For every inhabited pre-shift tag $k\in\mathcal N_{3}^{\mathrm{inh}}$,
\[
 \mathrm{B}_{4}(\Sigma_{4}(k))=k.
\]
For every inhabited shifted tag $a\geq 1$,
\[
 \Sigma_{4}(\mathrm{B}_{4}(a))=a.
\]
\end{lemma}

\begin{proof}
Both identities are the elementary identities $(k+1)-1=k$ and $(a-1)+1=a$.
The reserved value $0$ is excluded from the inhabited backtranslation domain by
definition. Thus the round-trip statement is exact on inhabited tags and
honestly undefined on the reserved empty marker.
\end{proof}

\subsection{Migration obligations}
\label{subsec:v3370-migration-obligations}

\begin{definition}[shifted migration obligation list]
\label{def:v3370-shifted-migration-obligation-list}
A later shifted global QAM Hilbert-shift must discharge the following obligations:
\begin{enumerate}[leftmargin=*]
 \item translate every inhabited QAM tag through $\Sigma_{4}$;
 \item reserve tag $0$ globally for the empty or no-active-kind marker;
 \item preserve all payloads attached to translated tags;
 \item translate route-mask examples and bitstream examples coherently;
 \item preserve the two-sided structural tensor grammar under tag erasure and
 reindexing;
 \item preserve router-object guards and backtranslation records;
 \item state explicitly that the HC-boundary activation shift is not being
 redefined by the global tag migration;
 \item prove a conservativity theorem by erasing shifted tags back to the
 pre-shift namespace through $\mathrm{B}_{4}$ on inhabited components.
\end{enumerate}
\end{definition}

\begin{remark}[What is not migrated]
\label{rem:v3370-what-is-not-migrated}
The shifted tag migration should not change classical upper--lower tensor
contraction, same-variance comparison rules, the OP status, the Historical
Witness, or the physical meaning of an already guarded readout. It changes
the QAM namespace and routing labels, not the underlying mathematical payloads.
\end{remark}

\subsection{Router-package migration}
\label{subsec:v3370-router-package-migration}

\begin{definition}[Router-package translation]
\label{def:v3370-router-package-translation}
Let
\[
 \mathfrak R=(K,M,C,T,B,G)
\]
be a guarded router object, with kind registry $K$, route-mask registry $M$,
classical collars $C$, trace record $T$, backtranslation record $B$, and guard
record $G$. Its shifted translation is the object
\[
 \Sigma_{4}(\mathfrak R)
 =
 (\Sigma_{4}K,\Sigma_{4}M,C,T,\mathrm{B}_{4}^{\mathfrak R},G),
\]
where collars, traces, and guards are preserved, while kind and mask registries
are reindexed and the backtranslation record is replaced by the corresponding
shifted-to-pre-shift decoding data.
\end{definition}

\begin{lemma}[Router guard preservation]
\label{lem:v3370-router-guard-preservation}
If a router object $\mathfrak R$ is guarded before migration, then
$\Sigma_{4}(\mathfrak R)$ carries the same guard record after migration.
\end{lemma}

\begin{proof}
By Definition~\ref{def:v3370-router-package-translation}, the guard record
$G$ is copied unchanged. The migration reindexes registries and replaces the
backtranslation ledger; it does not modify the guard predicates. Hence guard
status is preserved.
\end{proof}

\subsection{Readout-package migration}
\label{subsec:v3370-readout-package-migration}

\begin{definition}[Readout-package translation]
\label{def:v3370-readout-package-translation}
A guarded readout package may be translated to shifted only by translating its
QAM/router labels while leaving its declared readout target, comparison tensor,
boundary collar, residual map, and non-claim guard unchanged.
\end{definition}

\begin{proposition}[Machine/content preservation under migration]
\label{prop:v3370-machine-content-preservation}
The future shifted global QAM Hilbert-shift preserves the historical release line machine/content
separation: QAM/MML/MMA remains the router/compiler language, and physical or
mathematical content remains in the guarded readout package.
\end{proposition}

\begin{proof}
The migration acts on labels, kinds, masks, and routing registries. It does not
act on the declared target of a physical readout, the comparison tensor, the
boundary collar, or the residual map. Therefore it changes the machine
addressing language while preserving the separation between machine layer and
content layer.
\end{proof}

\subsection{Conservativity statement for the future shifted shift}
\label{subsec:v3370-conservativity-statement-future-shifted-shift}

\begin{theorem}[shifted migration conservativity contract]
\label{thm:v3370-shifted-migration-conservativity-contract}
Suppose a future shifted text translates every inhabited pre-shift QAM tag through
$\Sigma_{4}$, reserves $0$ only for the empty/no-active-kind marker, and equips
each translated router or readout object with its backtranslation ledger
$\mathrm{B}_{4}$. Then every finite guarded shifted router/readout package has an
inhabited-component backtranslation to a finite guarded pre-shift package with the
same payloads and the same theorematic status.
\end{theorem}

\begin{proof}
Apply $\mathrm{B}_{4}$ to every inhabited shifted tag and leave non-tag payloads fixed.
Round-trip correctness on inhabited tags follows from
Lemma~\ref{lem:v3370-round-trip-inhabited-tags}. Router guards are preserved
by Lemma~\ref{lem:v3370-router-guard-preservation}. Readout targets,
comparison tensors, boundary collars, residual maps, and non-claim guards are
unchanged by Definition~\ref{def:v3370-readout-package-translation}. Hence
the translated object backtranslates to a pre-shift guarded object with identical
payloads and unchanged theorematic status. The reserved value $0$ has no
inhabited backtranslation and therefore cannot be used to manufacture a new
certificate.
\end{proof}

\subsection{historical release line handoff to shifted}
\label{subsec:v3370-handoff-to-shifted}

\begin{definition}[shifted entry checklist]
\label{def:v3370-shifted-entry-checklist}
A later \QSCver{shifted.0.0} draft should begin only after the following checklist is
satisfied:
\begin{enumerate}[leftmargin=*]
 \item the global tag translation $\Sigma_{4}$ is fixed;
 \item every occurrence class of QAM tags has an entry in the migration ledger;
 \item every route-mask example has a shifted and pre-shift display pair;
 \item every router object carries a backtranslation record;
 \item the HC-boundary activation shift remains local and is not redefined;
 \item the shifted text carries the conservativity contract of
 Theorem~\ref{thm:v3370-shifted-migration-conservativity-contract};
 \item OP~1/4 and OP~5(B) are not claimed closed merely by the migration.
\end{enumerate}
\end{definition}

\begin{remark}[historical release line non-claim]
\label{rem:v3370-non-claim}
Version~3.37.0 is a migration-map and backtranslation-contract layer. It does
not execute the shifted global QAM Hilbert-shift, does not close OP~1/4 or OP~5(B),
does not activate primitive mixed tensors, and does not claim
Turing-completeness. Its purpose is to make the later shifted migration auditable
before the namespace is actually changed.
\end{remark}



\clearpage
\section{\QSCver{legacy.37.1} Concrete Migration Ledger and Route-Mask Examples}
\label{sec:v3371-concrete-migration-ledger-route-mask-examples}

\begin{quote}\small
\noindent\textbf{Status.}
This section is a concrete-audit layer for the historical release line migration map. It still
does not execute the global shifted QAM Hilbert-shift. Its purpose is to make the
future migration visually and mechanically checkable before the namespace is
changed.
\end{quote}

\subsection{Minimal tag ledger}
\label{subsec:v3371-minimal-tag-ledger}

\begin{definition}[Concrete migration ledger]
\label{def:v3371-concrete-migration-ledger}
The concrete migration ledger is the finite display of the rule
\[
 \Sigma_4(k)=k+1,
 \qquad
 \mathrm B_4(a)=a-1\quad(a\geq 1),
\]
together with the reserved value
\[
 0_{\mathrm{shift}}=\varnothing_{\mathrm{active}},
\]
which has no inhabited pre-shift backtranslation.
\end{definition}

\begin{center}
\small
\begin{tabular}{@{}llll@{}}
\toprule
pre-shift tag & shifted tag & reading & backtranslation\\
\midrule
reserved / none & $0$ & empty/no-active-kind & none\\
$0$ & $1$ & first inhabited kind & $\mathrm B_4(1)=0$\\
$1$ & $2$ & second inhabited kind & $\mathrm B_4(2)=1$\\
$2$ & $3$ & third inhabited kind & $\mathrm B_4(3)=2$\\
$k$ & $k+1$ & inhabited kind $k$ & $\mathrm B_4(k+1)=k$\\
\bottomrule
\end{tabular}
\end{center}

\begin{remark}[Why the ledger is not a new theorem]
\label{rem:v3371-ledger-not-new-theorem}
The ledger is a bookkeeping surface. It records how inhabited tags will be
renamed in shifted. It does not change the underlying readout, comparison tensor,
boundary collar, residual map, or OP status.
\end{remark}

\subsection{Route payload example}
\label{subsec:v3371-route-payload-example}

\begin{definition}[Route payload migration]
\label{def:v3371-route-payload-migration}
For a finite ordered route payload
\[
 r=\langle k_1,\ldots,k_n\rangle,
\]
the future shifted route payload is
\[
 \Sigma_4(r)
 =
 \langle k_1+1,\ldots,k_n+1\rangle.
\]
Its backtranslation is applied componentwise:
\[
 \mathrm B_4(\langle a_1,\ldots,a_n\rangle)
 =
 \langle a_1-1,\ldots,a_n-1\rangle
\]
whenever all $a_i\geq 1$.
\end{definition}

\begin{remark}[Example: two-entry route payload]
\label{rem:v3371-two-entry-route-payload-example}
The pre-shift payload
\[
 r=\langle 5,12\rangle
\]
migrates to
\[
 \Sigma_4(r)=\langle 6,13\rangle.
\]
Backtranslation returns
\[
 \mathrm B_4(\langle 6,13\rangle)=\langle 5,12\rangle.
\]
Thus the route addresses change, while the ordered route structure and the
payload role are preserved.
\end{remark}

\subsection{Packed mask example}
\label{subsec:v3371-packed-mask-example}

\begin{definition}[Support mask versus packed mask]
\label{def:v3371-support-mask-versus-packed-mask}
A support mask is a finite inhabited support set
\[
 S\subseteq \mathbb N.
\]
Its pre-shift packed integer is
\[
 m_3(S)=\sum_{k\in S}2^k.
\]
After migration the support becomes
\[
 \Sigma_4(S)=\{k+1:k\in S\},
\]
and its shifted packed integer is
\[
 m_4(\Sigma_4(S))=\sum_{k\in S}2^{k+1}.
\]
\end{definition}

\begin{lemma}[Packed-mask doubling rule]
\label{lem:v3371-packed-mask-doubling-rule}
For every finite inhabited support $S$,
\[
 m_4(\Sigma_4(S))=2\,m_3(S).
\]
\end{lemma}

\begin{proof}
By definition,
\[
 m_4(\Sigma_4(S))
 =
 \sum_{k\in S}2^{k+1}
 =
 2\sum_{k\in S}2^k
 =
 2\,m_3(S).
\]
The rule applies to the packed integer representation of the support. It is
not the same operation as adding $1$ to the packed integer.
\end{proof}

\begin{remark}[Example: packed mask for the route support $\{5,12\}$]
\label{rem:v3371-packed-mask-five-twelve-example}
For
\[
 S=\{5,12\},
\]
the pre-shift packed mask is
\[
 m_3(S)=2^5+2^{12}=32+4096=4128.
\]
The shifted support is
\[
 \Sigma_4(S)=\{6,13\},
\]
and the shifted packed mask is
\[
 m_4(\Sigma_4(S))=2^6+2^{13}=64+8192=8256.
\]
This equals $2\cdot 4128$. Backtranslation of the support gives
\[
 \mathrm B_4(\{6,13\})=\{5,12\}.
\]
\end{remark}

\begin{remark}[No accidental activation of the reserved bit]
\label{rem:v3371-no-accidental-reserved-bit-activation}
Since inhabited supports are shifted upward, the shifted bit $0$ is not produced by
migration. It remains reserved for the empty/no-active-kind marker. This is
the operational reason why the global shift is compatible with the
HC-boundary interpretation: the first active bit is no longer confused with the
empty marker.
\end{remark}

\subsection{Two-sided tensor ledger example}
\label{subsec:v3371-two-sided-tensor-ledger-example}

\begin{definition}[Tensor-index migration ledger]
\label{def:v3371-tensor-index-migration-ledger}
A two-sided structural tensor expression may carry QAM/router tags inside its
front-side or back-side index payloads. A future shifted migration applies
$\Sigma_4$ only to those tag payloads. It does not change front/back sidedness,
upper/lower variance, or the declared contraction/comparison rule.
\end{definition}

\begin{remark}[Example: front-side route migration]
\label{rem:v3371-front-side-route-migration-example}
A readout-facing tensor package with front-side route payload
\[
 {}^{\langle 5,12\rangle}{}_{\langle 2\rangle}
 \mathcal T^{i}{}_{j}
\]
migrates to
\[
 {}^{\langle 6,13\rangle}{}_{\langle 3\rangle}
 \mathcal T^{i}{}_{j}.
\]
The front-side payloads are reindexed. The back-side structural indices
$i,j$ are not changed in this example. The tensor still has the same
front/back and upper/lower grammar.
\end{remark}

\begin{remark}[Example: full payload migration]
\label{rem:v3371-full-payload-migration-example}
If both front-side and back-side QAM payloads are explicitly tagged, then
\[
 {}^{\langle 5,12\rangle}{}_{\langle 2\rangle}
 \mathcal T^{\langle 7\rangle}{}_{\langle 10\rangle}
\]
migrates to
\[
 {}^{\langle 6,13\rangle}{}_{\langle 3\rangle}
 \mathcal T^{\langle 8\rangle}{}_{\langle 11\rangle}.
\]
This is a label migration only. It does not turn an upper index into a lower
index, a front-side readout into a back-side action, or a comparison tensor into
a contraction tensor.
\end{remark}

\begin{proposition}[Tensor grammar preservation under the migration ledger]
\label{prop:v3371-tensor-grammar-preservation-under-ledger}
The concrete migration ledger preserves the two-sided structural tensor
grammar.
\end{proposition}

\begin{proof}
The migration acts on tag payloads by $k\mapsto k+1$. The syntactic positions
of the payloads are unchanged: front-side payloads remain front-side payloads,
back-side payloads remain back-side payloads, upper payloads remain upper
payloads, and lower payloads remain lower payloads. Therefore classical
upper--lower readout and explicit same-variance comparison requirements are
unchanged.
\end{proof}

\subsection{shifted readiness consequence}
\label{subsec:v3371-shifted-readiness-consequence}

\begin{theorem}[Concrete ledger supports the shifted conservativity contract]
\label{thm:v3371-concrete-ledger-supports-shifted-conservativity-contract}
The concrete tag ledger, route payload example, packed-mask rule, and
two-sided tensor ledger examples support the historical release line backtranslation contract:
every migrated finite example displayed in this section backtranslates to its
pre-shift source example on inhabited components.
\end{theorem}

\begin{proof}
The tag ledger backtranslates by $\mathrm B_4(a)=a-1$ for $a\geq 1$. Route
payloads backtranslate componentwise. Packed masks backtranslate at the level
of support sets, not by subtracting $1$ from the packed integer. Tensor
payloads backtranslate by applying the same rule to the tags carried inside
their index payloads, while preserving sidedness and variance. Hence each
finite example returns to its pre-shift source example on inhabited components.
\end{proof}

\begin{remark}[\QSCver{legacy.37.1} non-claim]
\label{rem:v3371-non-claim}
Version~3.37.1 still does not execute the global shifted QAM Hilbert-shift. It only
adds concrete ledger examples. The actual global migration remains a shifted task.
\end{remark}



\clearpage
\section{Global QAM Hilbert-Shift Activation Candidate}
\label{sec:v400-r01-global-qam-hilbert-shift-activation}

\begin{quote}\small
\noindent\textbf{Status.}
This is the first shifted activation release-final state. It executes the global QAM
namespace decision prepared in Versions~3.37.0--3.37.1: from this layer
forward, the active QAM language reserves tag $0$ for the empty/no-active-kind
marker and reads inhabited tags in the shifted shifted namespace. The earlier
historical release line HC-boundary shift remains local and structural; the present layer is the
global language shift.
\end{quote}

\subsection{Governing shifted namespace convention}
\label{subsec:v400-r01-governing-shifted-namespace-convention}

\begin{definition}[Global shifted QAM Hilbert-shift convention]
\label{def:v400-r01-global-shifted-qam-hilbert-shift-convention}
The governing shifted QAM namespace is obtained from the inhabited pre-shift namespace
by the shift
\[
 \Sigma_4(k)=k+1.
\]
The tag $0$ in the shifted namespace is reserved globally for the empty or
no-active-kind marker:
\[
 0_{\mathrm{shift}}=\varnothing_{\mathrm{active}}.
\]
It is not the image of an inhabited pre-shift tag.
\end{definition}

\begin{definition}[shifted backtranslation ledger]
\label{def:v400-r01-shifted-backtranslation-ledger}
The shifted backtranslation ledger is the partial map
\[
 \mathrm B_4(a)=a-1
\]
defined on inhabited shifted tags $a\geq 1$. The reserved tag $0$ has no inhabited
backtranslation.
\end{definition}

\begin{remark}[What is executed in]
\label{rem:v400-r01-what-is-executed}
The executed decision is the global namespace convention: shifted readout packages
use the shifted tag convention by default. The text does not pretend that the
Historical Witness has been rewritten. Historical material remains readable by
the backtranslation ledger.
\end{remark}

\subsection{Compatibility with the first HC-boundary shift}
\label{subsec:v400-r01-compatibility-first-hc-boundary-shift}

\begin{definition}[Shift A and Shift B in the shifted text]
\label{def:v400-r01-shift-a-shift-b-shifted-text}
The shifted text keeps the following distinction.
\begin{enumerate}[leftmargin=*]
 \item \emph{Shift A} is the local historical release line HC-boundary activation
 \[
 0_{\mathrm{mask}}\longrightarrow 1_{\mathrm{mask}}.
 \]
 It explains why the Heisenberg Compensator appears as a boundary/readout
 normalizer at the first active readout transition.

 \item \emph{Shift B} is the global shifted QAM Hilbert-shift
 \[
 k\longmapsto k+1
 \]
 applied to inhabited QAM tags.
\end{enumerate}
\end{definition}

\begin{proposition}[HC-boundary shift is not redefined by shifted]
\label{prop:v400-r01-hc-boundary-shift-not-redefined}
The global shifted QAM Hilbert-shift does not redefine the historical release line HC-boundary
activation principle.
\end{proposition}

\begin{proof}
Shift A concerns the structural boundary from inactive mask state to first
active readout state. Shift B concerns the global label namespace for
inhabited QAM tags. These acts have different domains: the first is a local
boundary-normalization reading, the second is a global language migration.
Therefore the shifted namespace convention preserves the HC-boundary reading rather
than replacing it.
\end{proof}

\subsection{Governing shifted ledger}
\label{subsec:v400-r01-governing-shifted-ledger}

\begin{center}
\small
\begin{tabular}{@{}llll@{}}
\toprule
pre-shift tag & shifted tag & shifted reading & backtranslation\\
\midrule
none & $0$ & empty/no-active-kind & none\\
$0$ & $1$ & first inhabited kind & $\mathrm B_4(1)=0$\\
$1$ & $2$ & second inhabited kind & $\mathrm B_4(2)=1$\\
$2$ & $3$ & third inhabited kind & $\mathrm B_4(3)=2$\\
$k$ & $k+1$ & inhabited kind $k$ & $\mathrm B_4(k+1)=k$\\
\bottomrule
\end{tabular}
\end{center}

\begin{lemma}[shifted inhabited round-trip]
\label{lem:v400-r01-shifted-inhabited-round-trip}
For every inhabited pre-shift tag $k$,
\[
 \mathrm B_4(\Sigma_4(k))=k.
\]
For every inhabited shifted tag $a\geq 1$,
\[
 \Sigma_4(\mathrm B_4(a))=a.
\]
\end{lemma}

\begin{proof}
Both claims are the identities $(k+1)-1=k$ and $(a-1)+1=a$. The reserved
value $0$ is excluded from the inhabited backtranslation domain, so it cannot
generate an inhabited source tag.
\end{proof}

\subsection{Route-mask and support migration}
\label{subsec:v400-r01-route-mask-support-migration}

\begin{definition}[shifted route payload]
\label{def:v400-r01-shifted-route-payload}
A pre-shift ordered route payload
\[
 r=\langle k_1,\ldots,k_n\rangle
\]
is read in the shifted language as
\[
 r^{(4)}=\Sigma_4(r)=\langle k_1+1,\ldots,k_n+1\rangle.
\]
\end{definition}

\begin{definition}[shifted support and packed mask]
\label{def:v400-r01-shifted-support-and-packed-mask}
For a finite inhabited support $S\subseteq\mathbb N$, the migrated shifted support is
\[
 S^{(4)}=\{k+1:k\in S\}.
\]
If the pre-shift packed integer is
\[
 m_3(S)=\sum_{k\in S}2^k,
\]
then the shifted packed integer is
\[
 m_4(S^{(4)})=\sum_{k\in S}2^{k+1}.
\]
\end{definition}

\begin{lemma}[shifted packed-mask doubling]
\label{lem:v400-r01-shifted-packed-mask-doubling}
For every finite inhabited support $S$,
\[
 m_4(S^{(4)})=2\,m_3(S).
\]
\end{lemma}

\begin{proof}
The proof is the direct calculation
\[
 m_4(S^{(4)})
 =
 \sum_{k\in S}2^{k+1}
 =
 2\sum_{k\in S}2^k
 =
 2\,m_3(S).
\]
\end{proof}

\begin{remark}[Example: route and mask]
\label{rem:v400-r01-example-route-and-mask}
The pre-shift route payload $\langle 5,12\rangle$ is read in shifted as
$\langle 6,13\rangle$. The support $\{5,12\}$ has packed value
$2^5+2^{12}=4128$ before migration and packed value
$2^6+2^{13}=8256$ after migration. The packed integer doubles; the support
backtranslates by subtracting one from inhabited support entries.
\end{remark}

\subsection{Two-sided tensor grammar under shifted}
\label{subsec:v400-r01-two-sided-tensor-grammar-under-shifted}

\begin{definition}[shifted tensor-payload migration]
\label{def:v400-r01-shifted-tensor-payload-migration}
In a two-sided structural tensor expression, the shifted shift acts only on QAM or
router tag payloads carried by index positions. It does not alter front/back
sidedness, upper/lower variance, classical upper--lower contraction, or the
requirement that same-variance comparisons use explicit comparison or selector
tensors.
\end{definition}

\begin{remark}[Example: front-side migration]
\label{rem:v400-r01-example-front-side-migration}
The front-side route display
\[
 {}^{\langle 5,12\rangle}{}_{\langle 2\rangle}\mathcal T^{i}{}_{j}
\]
is read in shifted as
\[
 {}^{\langle 6,13\rangle}{}_{\langle 3\rangle}\mathcal T^{i}{}_{j}.
\]
The expression changes its tag payloads, not its tensor grammar.
\end{remark}

\begin{remark}[Example: full payload migration]
\label{rem:v400-r01-example-full-payload-migration}
If both front-side and back-side payloads are tagged, then
\[
 {}^{\langle 5,12\rangle}{}_{\langle 2\rangle}
 \mathcal T^{\langle 7\rangle}{}_{\langle 10\rangle}
\]
is read in shifted as
\[
 {}^{\langle 6,13\rangle}{}_{\langle 3\rangle}
 \mathcal T^{\langle 8\rangle}{}_{\langle 11\rangle}.
\]
The front/back positions and the upper/lower variance are preserved.
\end{remark}

\begin{theorem}[shifted tensor-grammar preservation]
\label{thm:v400-r01-shifted-tensor-grammar-preservation}
The global shifted QAM Hilbert-shift preserves the two-sided structural tensor
grammar.
\end{theorem}

\begin{proof}
The shift is a function on tag payloads. It does not move an index from front
to back, from back to front, from upper to lower, or from lower to upper.
Therefore the grammar of sidedness and variance is unchanged. Since classical
upper--lower contraction and explicit same-variance comparison are defined by
position and variance rather than by the numerical tag name, those rules are
preserved.
\end{proof}

\subsection{shifted router and readout activation}
\label{subsec:v400-r01-router-readout-activation}

\begin{definition}[shifted router object]
\label{def:v400-r01-shifted-router-object}
A shifted router object is a guarded router object whose inhabited kind registry and
route-mask registry are read in the shifted shifted namespace, while its collars,
trace records, backtranslation ledger, and guard record remain explicit.
\end{definition}

\begin{definition}[shifted readout object]
\label{def:v400-r01-shifted-readout-object}
A shifted readout object is a guarded readout package whose QAM/router labels are
read in the shifted shifted namespace, while its readout target, comparison tensor,
boundary collar, residual map, and non-claim guard remain unchanged.
\end{definition}

\begin{theorem}[shifted activation conservativity]
\label{thm:v400-r01-shifted-activation-conservativity}
Every finite guarded shifted router/readout package admitted by
Definitions~\ref{def:v400-r01-shifted-router-object} and
\ref{def:v400-r01-shifted-readout-object} backtranslates on inhabited components to
a finite guarded pre-shift package with the same payloads and the same theorematic
status.
\end{theorem}

\begin{proof}
Apply $\mathrm B_4(a)=a-1$ to every inhabited shifted tag. By
Lemma~\ref{lem:v400-r01-shifted-inhabited-round-trip}, this recovers the
corresponding pre-shift tag on inhabited components. Collars, trace records,
comparison tensors, readout targets, boundary collars, residual maps, and
non-claim guards are kept fixed by definition. The reserved value $0$ has no
inhabited backtranslation, so it cannot create a new active component. Hence
the package backtranslates without changing payloads or theorematic status.
\end{proof}

\subsection{handoff}
\label{subsec:v400-r01-handoff}

\begin{remark}[non-claims]
\label{rem:v400-r01-non-claims}
This shifted activation release-final state performs the namespace decision, but it does not
close OP~1/4 or OP~5(B), does not activate primitive mixed tensors, does not
claim Turing-completeness, and does not replace standard QFT/QED/Lifshitz or
general-relativistic calculations. It is a language and router/readout
migration step.
\end{remark}

\begin{theorem}[handoff theorem]
\label{thm:v400-r01-handoff-theorem}
The local proof line layer is a well-formed global QAM Hilbert-shift activation
release-final state over the \QSCver{legacy.37.1} release line.
\end{theorem}

\begin{proof}
Definition~\ref{def:v400-r01-global-shifted-qam-hilbert-shift-convention} fixes the
global namespace convention. Definition~\ref{def:v400-r01-shifted-backtranslation-ledger}
provides the backtranslation ledger. The route and mask rules preserve finite
payloads by Lemma~\ref{lem:v400-r01-shifted-packed-mask-doubling}. The two-sided
tensor grammar is preserved by Theorem~\ref{thm:v400-r01-shifted-tensor-grammar-preservation}.
Finally, finite router/readout packages remain conservative under
Theorem~\ref{thm:v400-r01-shifted-activation-conservativity}. Therefore the layer
is well formed as a global QAM language activation release-final state, while all stated
non-claims remain in force.
\end{proof}



\clearpage
\section{Occurrence Ledger and Shift-Audit Map}
\label{sec:v400-r02-occurrence-ledger-shift-audit-map}

\begin{quote}\small
\noindent\textbf{Status.}
This layer is an occurrence ledger for the shifted global QAM Hilbert-shift.
Version~4.0.0-fixed the governing shifted namespace. The present layer
records where that shift is allowed to act, where it is forbidden to act, and
where a backtranslation or historical reading must be retained. It is an audit
map, not a new physical promotion.
\end{quote}

\subsection{Occurrence classes}
\label{subsec:v400-r02-occurrence-classes}

\begin{definition}[shifted occurrence class]
\label{def:v400-r02-occurrence-class}
A \emph{shifted occurrence class} is a declared role assigned to a visible tag,
mask, route, tensor-index payload, or OP-facing label in the Active Core. The
role determines whether the global shifted shift is applied, whether the object is
kept historical, or whether the object is read only through an explicit
backtranslation ledger.
\end{definition}

\begin{definition}[The five primary occurrence classes]
\label{def:v400-r02-five-primary-occurrence-classes}
The occurrence ledger uses the following primary classes.
\begin{enumerate}[leftmargin=*]
 \item \textbf{Class A: historical/pre-shift witness occurrences.}
 These remain unchanged and are read through the Historical Witness or through
 the backtranslation ledger.

 \item \textbf{Class B: active shifted router occurrences.}
 These are read in the shifted shifted namespace by default.

 \item \textbf{Class C: example, table, route, and mask occurrences.}
 These must carry a shifted display and, when useful, a pre-shift backtranslation
 display.

 \item \textbf{Class D: two-sided tensor-index payload occurrences.}
 These shift only their QAM/router tag payloads. The tensor grammar itself is
 not shifted.

 \item \textbf{Class E: OP-facing and physical-readout occurrences.}
 These may have shifted router labels, but their theorematic status, physical
 readout target, residual map, and non-claim guard remain unchanged.
\end{enumerate}
\end{definition}

\begin{remark}[Why occurrence classes are needed]
\label{rem:v400-r02-why-occurrence-classes-needed}
Without occurrence classes, the symbol $5$ could ambiguously mean a pre-shift tag,
a shifted tag, a packed mask integer, a tensor component, or a mere example number.
The occurrence ledger prevents that ambiguity by assigning a role before any
shift is read.
\end{remark}

\subsection{Class A: historical and pre-shift witness occurrences}
\label{subsec:v400-r02-class-a-historical-pre-shift}

\begin{definition}[Class A occurrence]
\label{def:v400-r02-class-a-occurrence}
A Class A occurrence is any tag, route, mask, or notation block that belongs to
the preserved historical line, to the pre-shift witness text, or to an explicitly
quoted earlier version. Class A occurrences are not rewritten by the shifted Active
Core. They remain recoverable by the Historical Witness and readable by
backtranslation.
\end{definition}

\begin{lemma}[Class A preservation]
\label{lem:v400-r02-class-a-preservation}
The shifted occurrence ledger preserves Class A occurrences without applying the
global shift to their displayed historical form.
\end{lemma}

\begin{proof}
Class A is defined as historical or pre-shift witness material. Its function is
to preserve genealogy and reconstruction memory, not to govern the active shifted
namespace. Therefore applying the shifted shift directly to Class A would corrupt
the witness role. The correct operation is preservation plus optional
backtranslation annotation.
\end{proof}

\subsection{Class B: active shifted router occurrences}
\label{subsec:v400-r02-class-b-active-router}

\begin{definition}[Class B occurrence]
\label{def:v400-r02-class-b-occurrence}
A Class B occurrence is a tag or route occurrence inside an active shifted router
object, shifted kind registry, shifted route-mask registry, or shifted machine-language
package. Class B occurrences are read in the shifted namespace:
\[
 k_{\mathrm{pre}}\longmapsto k_{\mathrm{shift}}=k_{\mathrm{pre}}+1.
\]
\end{definition}

\begin{remark}[Class B default]
\label{rem:v400-r02-class-b-default}
From the local proof line onward, active router-language statements default to Class B
unless the text explicitly marks the occurrence as historical, illustrative, or
unshifted for comparison.
\end{remark}

\begin{lemma}[Class B backtranslation]
\label{lem:v400-r02-class-b-backtranslation}
Every finite Class B tag family with inhabited shifted tags has a componentwise
backtranslation by $\mathrm B_4(a)=a-1$.
\end{lemma}

\begin{proof}
Class B contains active inhabited shifted router labels. By the shifted ledger, every
inhabited shifted tag satisfies $a\geq 1$, so $\mathrm B_4(a)=a-1$ is defined.
Applying this componentwise gives the pre-shift label family.
\end{proof}

\subsection{Class C: examples, tables, routes, and masks}
\label{subsec:v400-r02-class-c-examples-tables-routes-masks}

\begin{definition}[Class C occurrence]
\label{def:v400-r02-class-c-occurrence}
A Class C occurrence is a displayed example, table entry, route payload, support
set, or packed mask integer used for exposition or audit. Class C occurrences
must declare whether the displayed value is a pre-shift value, a shifted value, or a
packed integer derived from a support.
\end{definition}

\begin{definition}[Class C display contract]
\label{def:v400-r02-class-c-display-contract}
A Class C route or mask example is well displayed if it separates:
\[
 \text{support},\qquad \text{route payload},\qquad \text{packed integer}.
\]
In particular, support migration is
\[
 S\mapsto S^{(4)}=\{k+1:k\in S\},
\]
while packed-mask migration is
\[
 m_3(S)=\sum_{k\in S}2^k
 \quad\mapsto\quad
 m_4(S^{(4)})=2m_3(S).
\]
\end{definition}

\begin{remark}[Class C warning]
\label{rem:v400-r02-class-c-warning}
The packed integer should not be shifted by adding $1$. The support is shifted
by adding $1$ to each inhabited support entry; the packed integer then changes
according to the support rule, often by doubling in the pure support-shift
case.
\end{remark}

\subsection{Class D: two-sided tensor-index payloads}
\label{subsec:v400-r02-class-d-two-sided-tensor-index-payloads}

\begin{definition}[Class D occurrence]
\label{def:v400-r02-class-d-occurrence}
A Class D occurrence is a QAM/router tag payload carried inside a front-side or
back-side index position of a two-sided structural tensor expression. The shifted
shift acts on the payload label only. It does not alter sidedness, variance,
ordinary upper--lower contraction, or explicit same-variance comparison.
\end{definition}

\begin{proposition}[Class D grammar invariance]
\label{prop:v400-r02-class-d-grammar-invariance}
Class D migration preserves the two-sided structural tensor grammar.
\end{proposition}

\begin{proof}
The migration rule changes only the numerical or route-label payload carried by
an index. It does not move that payload from the front side to the back side,
from the back side to the front side, from upper position to lower position, or
from lower position to upper position. Thus the grammar of sidedness and
variance is unchanged.
\end{proof}

\subsection{Class E: OP-facing and physical-readout occurrences}
\label{subsec:v400-r02-class-e-op-facing-physical-readout}

\begin{definition}[Class E occurrence]
\label{def:v400-r02-class-e-occurrence}
A Class E occurrence is an OP-facing label, residual package, physical-readout
package, boundary-readout package, or diagnostic object whose router address may
be shifted but whose physical or theorematic status is not changed by the
shift.
\end{definition}

\begin{definition}[Class E invariance contract]
\label{def:v400-r02-class-e-invariance-contract}
A Class E occurrence remains valid under shifted only if the following data remain
unchanged:
\begin{enumerate}[leftmargin=*]
 \item its OP status;
 \item its physical or mathematical readout target;
 \item its comparison tensor or boundary collar;
 \item its residual map or diagnostic map;
 \item its non-claim guard.
\end{enumerate}
\end{definition}

\begin{lemma}[Class E status preservation]
\label{lem:v400-r02-class-e-status-preservation}
The shifted occurrence ledger does not close OP~1/4 or OP~5(B).
\end{lemma}

\begin{proof}
Class E migration may reindex router labels, but
Definition~\ref{def:v400-r02-class-e-invariance-contract} fixes the OP status
and non-claim guard. Since no closure certificate is added by reindexing, the
OP status remains unchanged.
\end{proof}

\subsection{Occurrence decision table}
\label{subsec:v400-r02-occurrence-decision-table}

\begin{center}
\tiny
\setlength{\tabcolsep}{2pt}
\begin{tabular}{@{}p{0.10\textwidth}P{0.24\textwidth}P{0.23\textwidth}P{0.26\textwidth}@{}}
\toprule
Class & Where it occurs & Shift action & Required audit note\\
\midrule
A & Historical/pre-shift witness blocks & preserve displayed form & read via witness/backtranslation\\
B & active router/kind/route registries & shift inhabited tags & keep $\mathrm B_4$ ledger\\
C & examples, tables, masks, routes & display pre-shift and shifted roles & separate support from packed integer\\
D & tensor-index tag payloads & shift payload only & preserve sidedness/variance\\
E & OP-facing/readout packages & shift router address only & preserve OP status/guard\\
\bottomrule
\end{tabular}
\end{center}

\subsection{Ambiguity filters}
\label{subsec:v400-r02-ambiguity-filters}

\begin{definition}[Ambiguity filter]
\label{def:v400-r02-ambiguity-filter}
Before a visible integer, tuple, support, or tensor payload is read as shifted,
the reader must determine whether it is:
\[
\begin{array}{lll}
 \text{tag},&
 \text{route payload},&
 \text{support set},\\
 \text{packed mask integer},&
 \text{tensor component},&
 \text{historical witness value}.
\end{array}
\]
Only tag, route-payload, support, and tag-payload occurrences are subject to
the shifted shift.
\end{definition}

\begin{remark}[Packed values are not tags]
\label{rem:v400-r02-packed-values-are-not-tags}
A packed integer such as $4128$ is not itself a tag occurrence. It encodes a
support. The support migrates first, and the packed integer is recomputed from
the migrated support.
\end{remark}

\subsection{shifted occurrence-ledger conservativity}
\label{subsec:v400-r02-occurrence-ledger-conservativity}

\begin{theorem}[Occurrence-ledger conservativity]
\label{thm:v400-r02-occurrence-ledger-conservativity}
Every finite shifted occurrence family classified by Classes A--E has an honest
reading that either preserves its displayed historical form or backtranslates
its inhabited shifted components to the pre-shift source language without changing
payloads, tensor grammar, OP status, or non-claim guards.
\end{theorem}

\begin{proof}
Class A is preserved by Lemma~\ref{lem:v400-r02-class-a-preservation}. Class B
backtranslates by Lemma~\ref{lem:v400-r02-class-b-backtranslation}. Class C
separates support, route payload, and packed integer by
Definition~\ref{def:v400-r02-class-c-display-contract}. Class D preserves the
two-sided tensor grammar by Proposition~\ref{prop:v400-r02-class-d-grammar-invariance}.
Class E preserves OP and readout status by Lemma~\ref{lem:v400-r02-class-e-status-preservation}.
The five cases cover the declared occurrence classes; hence the occurrence
ledger is conservative for finite classified occurrence families.
\end{proof}

\begin{remark}[non-claim]
\label{rem:v400-r02-non-claim}
The occurrence ledger does not add a new physical equation, does not close
an OP, does not activate primitive mixed tensors, and does not claim
Turing-completeness. It makes the global shifted shift auditable by assigning
roles to visible occurrences.
\end{remark}



\clearpage
\section{Global Reindexed Display Layer}
\label{sec:v400-r03-global-reindexed-display-layer}

\begin{quote}\small
\noindent\textbf{Status.}
This layer adds normalized side-by-side displays for the global shifted QAM
Hilbert-shift. It does not add new physics, close OP~1/4, close OP~5(B), or
alter the two-sided tensor grammar. It only standardizes how active examples
are displayed in three synchronized forms:
\[
 \text{pre-shift source},\qquad
 \text{shifted display},\qquad
 \text{backtranslation}.
\]
\end{quote}

\subsection{Reindexed display triple}
\label{subsec:v400-r03-reindexed-display-triple}

\begin{definition}[Reindexed display triple]
\label{def:v400-r03-reindexed-display-triple}
A \emph{shifted reindexed display triple} is a finite record
\[
 \mathfrak D_4(X)
 =
 (X_{\mathrm{pre}},X_{\mathrm{shift}},\mathrm B_4(X_{\mathrm{shift}})),
\]
where $X_{\mathrm{pre}}$ is the pre-shift source display, $X_{\mathrm{shift}}$ is the
shifted shifted display, and $\mathrm B_4(X_{\mathrm{shift}})$ is the explicit
backtranslation on inhabited components.
\end{definition}

\begin{remark}[Display triples are audit objects]
\label{rem:v400-r03-display-triples-are-audit-objects}
A display triple is not a new theorematic certificate. It is an audit object
showing that the visible shifted display can be read back to the pre-shift source where
the components are inhabited.
\end{remark}

\subsection{Tag and route display triples}
\label{subsec:v400-r03-tag-route-display-triples}

\begin{definition}[Tag display triple]
\label{def:v400-r03-tag-display-triple}
For an inhabited pre-shift tag $k$, the shifted tag display triple is
\[
 \mathfrak D_4(k)
 =
 (k,\ k+1,\ \mathrm B_4(k+1)=k).
\]
The reserved value $0_{\mathrm{shift}}$ has no inhabited pre-shift source and is
displayed separately as
\[
 (\varnothing_{\mathrm{active}},\ 0_{\mathrm{shift}},\ \text{no inhabited backtranslation}).
\]
\end{definition}

\begin{definition}[Route display triple]
\label{def:v400-r03-route-display-triple}
For a finite pre-shift route payload
\[
 r=\langle k_1,\ldots,k_n\rangle,
\]
the shifted display triple is
\[
 \mathfrak D_4(r)
 =
 \left(
 \langle k_1,\ldots,k_n\rangle,\ 
 \langle k_1+1,\ldots,k_n+1\rangle,\ 
 \langle k_1,\ldots,k_n\rangle
 \right).
\]
\end{definition}

\begin{center}
\small
\begin{tabular}{@{}lll@{}}
\toprule
pre-shift source & shifted display & backtranslation\\
\midrule
$5$ & $6$ & $\mathrm B_4(6)=5$\\
$12$ & $13$ & $\mathrm B_4(13)=12$\\
$\langle 5,12\rangle$ & $\langle 6,13\rangle$ & $\langle 5,12\rangle$\\
\bottomrule
\end{tabular}
\end{center}

\subsection{Support and packed-mask display triples}
\label{subsec:v400-r03-support-packed-mask-display-triples}

\begin{definition}[Support display triple]
\label{def:v400-r03-support-display-triple}
For a finite inhabited support set $S$, the support display triple is
\[
 \mathfrak D_4(S)
 =
 (S,\ S^{(4)},\ \mathrm B_4(S^{(4)})),
 \qquad
 S^{(4)}=\{k+1:k\in S\}.
\]
\end{definition}

\begin{definition}[Packed-mask display triple]
\label{def:v400-r03-packed-mask-display-triple}
For a finite inhabited support $S$, the packed-mask display triple is
\[
 \mathfrak D_4(m(S))
 =
 \left(
 m_3(S),\
 m_4(S^{(4)}),\
 S=\mathrm B_4(S^{(4)})
 \right).
\]
The backtranslation is attached to the support, not to the packed integer by
subtraction.
\end{definition}

\begin{center}
\small
\begin{tabular}{@{}lll@{}}
\toprule
pre-shift source & shifted display & backtranslation\\
\midrule
$\{5,12\}$ & $\{6,13\}$ & $\{5,12\}$\\
$2^5+2^{12}=4128$ & $2^6+2^{13}=8256$ & support returns to $\{5,12\}$\\
\bottomrule
\end{tabular}
\end{center}

\begin{lemma}[Packed-mask display consistency]
\label{lem:v400-r03-packed-mask-display-consistency}
The packed-mask display triple for a finite support $S$ is consistent with the
packed-mask doubling rule:
\[
 m_4(S^{(4)})=2\,m_3(S).
\]
\end{lemma}

\begin{proof}
This is exactly Lemma~\ref{lem:v400-r02-class-b-backtranslation} applied to the
inhabited support entries together with the packed-mask calculation
\[
 \sum_{k\in S}2^{k+1}=2\sum_{k\in S}2^k.
\]
The displayed packed integer is recomputed from the migrated support, not
shifted as a tag.
\end{proof}

\subsection{Two-sided tensor display triples}
\label{subsec:v400-r03-two-sided-tensor-display-triples}

\begin{definition}[Tensor display triple]
\label{def:v400-r03-tensor-display-triple}
A two-sided structural tensor display triple is obtained by applying the shifted
tag shift to QAM/router payloads inside index positions while leaving sidedness
and variance unchanged.
\end{definition}

\begin{center}
\tiny
\setlength{\tabcolsep}{1.5pt}
\begin{tabular}{@{}p{0.29\textwidth}P{0.29\textwidth}P{0.29\textwidth}@{}}
\toprule
pre-shift source & shifted display & backtranslation\\
\midrule
${}^{\langle 5,12\rangle}{}_{\langle 2\rangle}\mathcal T^{i}{}_{j}$
&
${}^{\langle 6,13\rangle}{}_{\langle 3\rangle}\mathcal T^{i}{}_{j}$
&
${}^{\langle 5,12\rangle}{}_{\langle 2\rangle}\mathcal T^{i}{}_{j}$
\\[0.6em]
${}^{\langle 5,12\rangle}{}_{\langle 2\rangle}\mathcal T^{\langle 7\rangle}{}_{\langle 10\rangle}$
&
${}^{\langle 6,13\rangle}{}_{\langle 3\rangle}\mathcal T^{\langle 8\rangle}{}_{\langle 11\rangle}$
&
${}^{\langle 5,12\rangle}{}_{\langle 2\rangle}\mathcal T^{\langle 7\rangle}{}_{\langle 10\rangle}$
\\
\bottomrule
\end{tabular}
\end{center}

\begin{proposition}[Display triples preserve tensor grammar]
\label{prop:v400-r03-display-triples-preserve-tensor-grammar}
Tensor display triples preserve the two-sided structural tensor grammar.
\end{proposition}

\begin{proof}
Only QAM/router tag payloads inside the displayed index positions are shifted.
The index positions themselves are not moved. Hence front-side remains
front-side, back-side remains back-side, upper remains upper, lower remains
lower, and the ordinary contraction/comparison rules remain those of the
underlying tensor grammar.
\end{proof}

\subsection{OP-facing display triples}
\label{subsec:v400-r03-op-facing-display-triples}

\begin{definition}[OP-facing display triple]
\label{def:v400-r03-op-facing-display-triple}
An OP-facing display triple records only the router-address migration of an
OP-facing package. Its OP status, residual target, comparison surface, and
non-claim guard are displayed as invariant data.
\end{definition}

\begin{center}
\scriptsize
\setlength{\tabcolsep}{2pt}
\begin{tabular}{@{}p{0.24\textwidth}P{0.25\textwidth}P{0.35\textwidth}@{}}
\toprule
pre-shift address & shifted address & invariant data\\
\midrule
OP package tag $k$ & OP package tag $k+1$ & OP status unchanged; guard unchanged\\
residual route $\langle 5,12\rangle$ & residual route $\langle 6,13\rangle$ & residual target unchanged\\
boundary readout route $\langle 2\rangle$ & boundary readout route $\langle 3\rangle$ & boundary collar unchanged\\
\bottomrule
\end{tabular}
\end{center}

\begin{lemma}[OP-facing display triples do not close OPs]
\label{lem:v400-r03-op-facing-display-triples-do-not-close-ops}
OP-facing display triples do not close OP~1/4 or OP~5(B).
\end{lemma}

\begin{proof}
By Definition~\ref{def:v400-r03-op-facing-display-triple}, the triple changes
only router-address display. It explicitly leaves OP status, residual target,
comparison surface, and non-claim guard invariant. Therefore no OP closure
certificate is introduced.
\end{proof}

\subsection{Display-layer conservativity}
\label{subsec:v400-r03-display-layer-conservativity}

\begin{theorem}[Global reindexed display conservativity]
\label{thm:v400-r03-global-reindexed-display-conservativity}
Every finite display triple has a backtranslation to its pre-shift source
display on inhabited components and preserves the occurrence class assigned in
the occurrence ledger.
\end{theorem}

\begin{proof}
Tag and route triples backtranslate componentwise by
$\mathrm B_4(a)=a-1$. Support triples backtranslate by applying the same rule
to support entries. Packed-mask triples backtranslate through their support
sets, not by subtracting from the packed integer. Tensor triples preserve
grammar by Proposition~\ref{prop:v400-r03-display-triples-preserve-tensor-grammar}.
OP-facing triples preserve status by
Lemma~\ref{lem:v400-r03-op-facing-display-triples-do-not-close-ops}. Hence the
display layer is conservative over the occurrence ledger.
\end{proof}

\begin{remark}[non-claim]
\label{rem:v400-r03-non-claim}
The layer is a display layer only. It does not add a new physical equation,
does not close OP~1/4 or OP~5(B), does not activate primitive mixed tensors, and
does not claim Turing-completeness.
\end{remark}



\clearpage
\section{Physical Readout Compatibility Check}
\label{sec:v400-r04-physical-readout-compatibility-check}

\begin{quote}\small
\noindent\textbf{Status.}
This layer checks that the global shifted QAM Hilbert-shift is compatible with
the active physical and mathematical readout packages. It does not promote a
new physical law. It verifies that the shifted shift changes router addresses and
tag payloads only, while readout targets, comparison surfaces, residual maps,
and non-claim guards remain fixed.
\end{quote}

\subsection{Readout compatibility principle}
\label{subsec:v400-r04-readout-compatibility-principle}

\begin{definition}[Readout compatibility package]
\label{def:v400-r04-readout-compatibility-package}
A \emph{shifted readout compatibility package} is a tuple
\[
 \mathcal C_4(R)
 =
 (R_{\mathrm{pre}},R_{\mathrm{shift}},\mathrm B_4(R_{\mathrm{shift}}),
 T_R,G_R,S_R),
\]
where $R_{\mathrm{pre}}$ is the pre-shift router address, $R_{\mathrm{shift}}$ is the
shifted shifted router address, $\mathrm B_4(R_{\mathrm{shift}})$ is the inhabited
backtranslation, $T_R$ is the readout target, $G_R$ is the guard record, and
$S_R$ is the theorematic status record.
\end{definition}

\begin{definition}[Readout compatibility]
\label{def:v400-r04-readout-compatibility}
A readout package is \emph{shifted-compatible} if the migration from
$R_{\mathrm{pre}}$ to $R_{\mathrm{shift}}$ changes only the router address and
leaves $T_R$, $G_R$, and $S_R$ unchanged.
\end{definition}

\begin{lemma}[Readout compatibility is address-local]
\label{lem:v400-r04-readout-compatibility-address-local}
If a readout package is shifted-compatible, then the global shifted QAM Hilbert-shift
does not change its physical or mathematical target.
\end{lemma}

\begin{proof}
By Definition~\ref{def:v400-r04-readout-compatibility}, the shift is applied to
the router address only. The readout target $T_R$ is explicitly preserved.
Hence the physical or mathematical target of the package is unchanged.
\end{proof}

\subsection{HC-boundary readout compatibility}
\label{subsec:v400-r04-hc-boundary-readout-compatibility}

\begin{definition}[HC-boundary compatibility record]
\label{def:v400-r04-hc-boundary-compatibility-record}
The HC-boundary readout package is shifted-compatible when the local structural
shift
\[
 0_{\mathrm{mask}}\longrightarrow 1_{\mathrm{mask}}
\]
remains the HC-boundary activation principle, while global shifted router labels are
read in the shifted namespace.
\end{definition}

\begin{proposition}[HC-boundary readout survives the global shift]
\label{prop:v400-r04-hc-boundary-survives-global-shift}
The global shifted QAM Hilbert-shift preserves the historical release line HC-boundary activation
reading.
\end{proposition}

\begin{proof}
The HC-boundary reading belongs to Shift A: the local transition from inactive
mask state to first active readout state. The shifted QAM Hilbert-shift belongs to
Shift B: the global namespace migration of inhabited tags. Since the local proof chain
separate these two acts, the global shift cannot redefine the local HC-boundary
activation. It can only change how the active router address is displayed.
\end{proof}

\subsection{Photon-compatible null-side readout compatibility}
\label{subsec:v400-r04-photon-null-side-readout-compatibility}

\begin{definition}[Null-side readout compatibility]
\label{def:v400-r04-null-side-readout-compatibility}
A photon-compatible null-side readout is shifted-compatible when the null-side
readout target and its guard remain unchanged while its router address is
displayed in the shifted shifted namespace.
\end{definition}

\begin{lemma}[Null-side target preservation]
\label{lem:v400-r04-null-side-target-preservation}
The shifted shift preserves the null-side readout target of a photon-compatible
readout package.
\end{lemma}

\begin{proof}
By Definition~\ref{def:v400-r04-null-side-readout-compatibility}, the null-side
target is part of the invariant data. The shifted shift operates only on router
address labels, so the null-side target remains fixed.
\end{proof}

\begin{remark}[No photon proof claim]
\label{rem:v400-r04-no-photon-proof-claim}
This compatibility check does not prove a new photon theory. It only says that
the already guarded null-side readout remains address-stable under the shifted tag
migration.
\end{remark}

\subsection{HC-normalized singularity projection compatibility}
\label{subsec:v400-r04-singularity-projection-compatibility}

\begin{definition}[Singularity projection compatibility]
\label{def:v400-r04-singularity-projection-compatibility}
An HC-normalized singularity projection package is shifted-compatible when its
projection target, HC-normalization guard, and residual diagnostic remain
unchanged under router-address migration.
\end{definition}

\begin{proposition}[Singularity projection status preservation]
\label{prop:v400-r04-singularity-projection-status-preservation}
The shifted shift does not promote an HC-normalized singularity projection package
from diagnostic status to theorematic closure.
\end{proposition}

\begin{proof}
The migration changes the router address. It does not add a new projection
identity, a new residual estimate, or a new closure certificate. The
diagnostic guard is part of the invariant record in
Definition~\ref{def:v400-r04-singularity-projection-compatibility}. Therefore
the theorematic status is preserved.
\end{proof}

\subsection{Casimir boundary readout compatibility}
\label{subsec:v400-r04-casimir-boundary-readout-compatibility}

\begin{definition}[Casimir boundary compatibility]
\label{def:v400-r04-casimir-boundary-compatibility}
A Casimir boundary readout package is shifted-compatible when its boundary collar,
comparison surface, and non-replacement guard remain unchanged while its
router-address labels are shifted.
\end{definition}

\begin{lemma}[Casimir boundary collar preservation]
\label{lem:v400-r04-casimir-boundary-collar-preservation}
The shifted shift preserves the boundary collar of a shifted-compatible Casimir boundary
readout package.
\end{lemma}

\begin{proof}
By Definition~\ref{def:v400-r04-casimir-boundary-compatibility}, the boundary
collar is invariant data. The shift acts only on router labels. Therefore the
boundary collar is unchanged.
\end{proof}

\begin{remark}[No replacement of standard Casimir calculations]
\label{rem:v400-r04-no-replacement-standard-casimir}
The compatibility statement does not replace standard QFT, QED, or Lifshitz
calculations. It keeps the Companion readout package address-stable under the
shifted namespace shift.
\end{remark}

\subsection{OP-facing residual package compatibility}
\label{subsec:v400-r04-op-facing-residual-package-compatibility}

\begin{definition}[OP-facing residual compatibility]
\label{def:v400-r04-op-facing-residual-compatibility}
An OP-facing residual package is shifted-compatible when its route address may be
shifted but its residual target, status flag, and non-claim guard remain
unchanged.
\end{definition}

\begin{lemma}[OP-facing residual status preservation]
\label{lem:v400-r04-op-facing-residual-status-preservation}
The shifted shift does not close OP~1/4 or OP~5(B).
\end{lemma}

\begin{proof}
By Definition~\ref{def:v400-r04-op-facing-residual-compatibility}, OP-facing
residual packages preserve their status flags and non-claim guards. A router
address migration is not a proof certificate. Therefore the OP status remains
unchanged.
\end{proof}

\subsection{Compatibility table}
\label{subsec:v400-r04-compatibility-table}

\begin{center}
\scriptsize
\setlength{\tabcolsep}{2pt}
\begin{tabular}{@{}p{0.22\textwidth}P{0.30\textwidth}P{0.34\textwidth}@{}}
\toprule
Readout package & Shifted data & Invariant data\\
\midrule
HC-boundary & active router address & local HC-boundary activation principle\\
null-side/photon-compatible & router address & null-side target and guard\\
singularity projection & router address & projection target, HC guard, residual diagnostic\\
Casimir boundary & router address & boundary collar, comparison surface, non-replacement guard\\
OP-facing residual & route address & residual target, OP status, non-claim guard\\
\bottomrule
\end{tabular}
\end{center}

\subsection{Physical-readout compatibility theorem}
\label{subsec:v400-r04-physical-readout-compatibility-theorem}

\begin{theorem}[Physical readout compatibility under the shifted shift]
\label{thm:v400-r04-physical-readout-compatibility}
For the readout packages listed in Section~\ref{subsec:v400-r04-compatibility-table},
the global shifted QAM Hilbert-shift preserves the declared readout target,
comparison or boundary data, residual data, and non-claim guard. It changes
only the router-address display.
\end{theorem}

\begin{proof}
Each listed package is declared shifted-compatible by its corresponding definition:
HC-boundary compatibility by
Definition~\ref{def:v400-r04-hc-boundary-compatibility-record}, null-side
compatibility by Definition~\ref{def:v400-r04-null-side-readout-compatibility},
singularity projection compatibility by
Definition~\ref{def:v400-r04-singularity-projection-compatibility}, Casimir
boundary compatibility by
Definition~\ref{def:v400-r04-casimir-boundary-compatibility}, and OP-facing
residual compatibility by
Definition~\ref{def:v400-r04-op-facing-residual-compatibility}. In each case,
the shifted component is the router or route address. The physical or
mathematical readout target and the guard/status records are invariant data.
Therefore the shifted shift is compatible with these readout packages without
changing their theorematic status.
\end{proof}

\begin{remark}[non-claim]
\label{rem:v400-r04-non-claim}
The layer is a compatibility check. It does not add a new physical
equation, does not close OP~1/4 or OP~5(B), does not activate primitive mixed
tensors, and does not claim Turing-completeness.
\end{remark}



\clearpage
\section{Governing shifted Normal Form and Release-Freeze Handoff}
\label{sec:v400-r05-governing-shifted-normal-form-release-freeze-handoff}

\begin{quote}\small
\noindent\textbf{Status.}
This layer compresses the \QSCver{shifted.0.0}-the local proof chain migration work into one governing
normal form. It is intended as a release-freeze handoff: after this layer,
new active QAM/router/readout material should be written either in the shifted
normal form or with an explicit pre-shift backtranslation record.
\end{quote}

\subsection{Governing normal-form object}
\label{subsec:v400-r05-governing-normal-form-object}

\begin{definition}[Governing shifted normal-form object]
\label{def:v400-r05-governing-shifted-normal-form-object}
A governing shifted normal-form object is a finite record
\[
 \mathcal N_4(X)
 =
 (C_X,\ X_{\mathrm{shift}},\ \mathrm B_4(X_{\mathrm{shift}}),\
 P_X,\ G_X,\ S_X),
\]
where $C_X$ is the occurrence class, $X_{\mathrm{shift}}$ is the shifted shifted
display, $\mathrm B_4(X_{\mathrm{shift}})$ is the inhabited-component
backtranslation, $P_X$ is the payload or readout target, $G_X$ is the guard
record, and $S_X$ is the theorematic status record.
\end{definition}

\begin{remark}[Why the pre-shift source is not always the first field]
\label{rem:v400-r05-pre-shift-source-not-always-first-field}
The shifted Active Core now governs by the shifted namespace. Therefore the normal
form puts the shifted display first and keeps the pre-shift source as a backtranslation
ledger when needed. Historical source displays remain available through the
Historical Witness.
\end{remark}

\subsection{Normal-form admissibility}
\label{subsec:v400-r05-normal-form-admissibility}

\begin{definition}[shifted-normal admissibility]
\label{def:v400-r05-shifted-normal-admissibility}
A finite object is \emph{shifted-normal admissible} if:
\begin{enumerate}[leftmargin=*]
 \item it carries an occurrence class $C_X$;
 \item all active inhabited tags are read in the shifted shifted namespace;
 \item every shifted inhabited tag has a backtranslation by $\mathrm B_4$;
 \item packed masks are reconstructed from shifted supports, not shifted as
 integers;
 \item tensor-index payload migration preserves sidedness and variance;
 \item physical/readout payloads preserve their target, collar, residual map,
 guard, and theorematic status.
\end{enumerate}
\end{definition}

\begin{lemma}[Normal admissibility implies ledger readability]
\label{lem:v400-r05-normal-admissibility-implies-ledger-readability}
Every shifted-normal admissible object is readable through the namespace ledger,
the occurrence ledger, the display triple, and the readout
compatibility check.
\end{lemma}

\begin{proof}
The occurrence class invokes the local pass. The shifted tag convention invokes the local pass. The
backtranslation field invokes whenever a visible pre-shift comparison is
needed. The payload/guard/status fields invoke for readout packages and
the tensor-grammar rules for tensor payloads. Hence the object is readable
through the whole the local proof chain ledger chain.
\end{proof}

\subsection{Normal forms by class}
\label{subsec:v400-r05-normal-forms-by-class}

\begin{definition}[Class-wise shifted-namespace normal forms]
\label{def:v400-r05-class-wise-shifted-normal-forms}
The following class-wise normal forms are used:
\[
\begin{array}{ll}
\text{Class A:} & (A,\ X_{\mathrm{hist}},\ \mathrm{witness/backtranslation},\ P_X,\ G_X,\ S_X),\\[0.2em]
\text{Class B:} & (B,\ X_{\mathrm{shift}},\ \mathrm B_4(X_{\mathrm{shift}}),\ P_X,\ G_X,\ S_X),\\[0.2em]
\text{Class C:} & (C,\ \text{shifted display},\ \text{pre-shift display},\ \text{support/payload split},\ G_X,\ S_X),\\[0.2em]
\text{Class D:} & (D,\ \text{shifted tag payload},\ \text{backtranslated payload},\ \text{tensor grammar},\ G_X,\ S_X),\\[0.2em]
\text{Class E:} & (E,\ \text{shifted route address},\ \text{backtranslated address},\ \text{readout target},\ G_X,\ S_X).
\end{array}
\]
\end{definition}

\begin{remark}[Class C special rule]
\label{rem:v400-r05-class-c-special-rule}
For Class C objects, a packed mask integer is never the primary shifted object.
The primary shifted object is the support or route payload. The packed integer
is recomputed from that shifted support.
\end{remark}

\subsection{Normal-form examples}
\label{subsec:v400-r05-normal-form-examples}

\begin{center}
\scriptsize
\setlength{\tabcolsep}{2pt}
\begin{tabular}{@{}p{0.14\textwidth}P{0.24\textwidth}P{0.25\textwidth}P{0.25\textwidth}@{}}
\toprule
Class & shifted display & backtranslation & invariant data\\
\midrule
B & $\langle 6,13\rangle$ & $\langle 5,12\rangle$ & router role\\
C & support $\{6,13\}$, mask $8256$ & support $\{5,12\}$, mask $4128$ & support/mask split\\
D & ${}^{\langle 6,13\rangle}{}_{\langle 3\rangle}\mathcal T^{i}{}_{j}$ & ${}^{\langle 5,12\rangle}{}_{\langle 2\rangle}\mathcal T^{i}{}_{j}$ & sidedness/variance\\
E & route $\langle 6,13\rangle$ & route $\langle 5,12\rangle$ & readout target and guard\\
\bottomrule
\end{tabular}
\end{center}

\subsection{Release-freeze discipline}
\label{subsec:v400-r05-release-freeze-discipline}

\begin{definition}[shifted release-freeze discipline]
\label{def:v400-r05-shifted-release-freeze-discipline}
The shifted release-freeze discipline is the rule that, after the local pass, every new active
QAM/router/readout object must either:
\begin{enumerate}[leftmargin=*]
 \item be written directly as a shifted-normal admissible object; or
 \item be marked as historical/pre-shift and supplied with an explicit
 backtranslation or witness pointer.
\end{enumerate}
\end{definition}

\begin{lemma}[Release-freeze prevents mixed namespace drift]
\label{lem:v400-r05-release-freeze-prevents-mixed-namespace-drift}
The shifted release-freeze discipline prevents uncontrolled mixing of pre-shift and shifted
tag namespaces.
\end{lemma}

\begin{proof}
Every new object must either be shifted-normal or explicitly historical/pre-shift. In
the first case, inhabited tags are shifted and backtranslated by the shifted ledger.
In the second case, the object is not treated as an active shifted governing object.
Thus no object may silently mix pre-shift and shifted tag meanings.
\end{proof}

\subsection{Governing shifted normal-form theorem}
\label{subsec:v400-r05-governing-shifted-normal-form-theorem}

\begin{theorem}[Governing shifted normal-form theorem]
\label{thm:v400-r05-governing-shifted-normal-form-theorem}
The normal form is conservative over the local proof chain shifted migration chain:
every finite shifted-normal admissible object has a controlled reading in the
shifted shifted namespace and an inhabited-component backtranslation that preserves
payloads, tensor grammar, physical/readout targets, guards, and theorematic
status.
\end{theorem}

\begin{proof}
Let $\mathcal N_4(X)$ be shifted-normal admissible. By
Definition~\ref{def:v400-r05-shifted-normal-admissibility}, it carries an occurrence
class and a backtranslation ledger for all inhabited shifted tags. Class A is
handled by witness preservation. Class B is handled by router backtranslation.
Class C is handled by support/payload separation and packed-mask recomputation.
Class D is handled by preservation of sidedness and variance. Class E is
handled by preservation of readout target, guard, and theorematic status.
These are exactly the local proof chain conservation mechanisms. Hence the normal form
is conservative over the shifted migration chain.
\end{proof}

\subsection{\QSCver{shifted.0.0} release-final handoff}
\label{subsec:v400-r05-release-final-handoff}

\begin{definition}[\QSCver{shifted.0.0} release-final readiness]
\label{def:v400-r05-v400-release-final-readiness}
The \QSCver{shifted.0.0} line is ready for a release-final pass when:
\begin{enumerate}[leftmargin=*]
 \item the shifted namespace is fixed;
 \item the occurrence ledger is present;
 \item the reindexed display triples are present;
 \item physical/readout compatibility is present;
 \item the governing shifted-namespace normal form is present;
 \item all non-claims remain explicit.
\end{enumerate}
\end{definition}

\begin{remark}[non-claim]
\label{rem:v400-r05-non-claim}
The layer is a normal-form and release-freeze handoff. It does not close
OP~1/4 or OP~5(B), does not activate primitive mixed tensors, does not claim
Turing-completeness, and does not replace standard physical calculations.
\end{remark}



\clearpage
\section{\QSCver{shifted.0.0} Release-Final Closure Note}
\label{sec:v400-release-final-closure-note}

\begin{quote}\small
\noindent\textbf{Status.}
This closure note records the \QSCver{shifted.0.0} Active Core as the first shifted
release-final state built from the local proof chain migration chain. The release-final state
contains the global QAM Hilbert-shift activation, occurrence-ledger audit map,
reindexed display layer, physical-readout compatibility check, and governing
shifted-namespace normal form. It records historical DOI-bearing status for that layer; current \QSCver{shifted.1.0} DOI metadata is governed by the frontmatter Version and DOI record.
\end{quote}

\subsection{Release-release-final state contents}
\label{subsec:v400-release-final-contents}

\begin{definition}[\QSCver{shifted.0.0} release-final content]
\label{def:v400-release-final-content}
The \QSCver{shifted.0.13} historical source-hygiene layer consists of the following recorded layers:
\begin{enumerate}[leftmargin=*]
 \item the local pass: global QAM Hilbert-shift activation;
 \item the local pass: occurrence ledger and shift-audit map;
 \item the local pass: global reindexed display layer;
 \item the local pass: physical readout compatibility check;
 \item the local pass: governing shifted-namespace normal form and release-freeze handoff.
\end{enumerate}
\end{definition}

\begin{definition}[\QSCver{shifted.0.0} release-final governing rule]
\label{def:v400-release-final-governing-rule}
From this release-final state onward, active QAM/router/readout objects are governed by
the shifted-namespace normal form unless explicitly marked as historical or pre-shift witness
material. Inhabited active tags are read in the shifted namespace, and the
reserved tag $0_{\mathrm{shift}}$ denotes the empty/no-active-kind marker.
\end{definition}

\subsection{Release-release-final state conservativity}
\label{subsec:v400-release-final-conservativity}

\begin{theorem}[\QSCver{shifted.0.0} release-final conservativity]
\label{thm:v400-release-final-conservativity}
Every finite active object admitted by the \QSCver{shifted.0.0} release-final normal form
has an inhabited-component backtranslation that preserves its payload, tensor
grammar, readout target, guard record, and theorematic status.
\end{theorem}

\begin{proof}
The result is the direct combination of the namespace ledger, the local pass
occurrence ledger, the display triples, the physical-readout
compatibility check, and the governing normal-form theorem. The ledger
supplies $\mathrm B_4(a)=a-1$ on inhabited shifted tags. The ledger assigns a
role to visible occurrences. The display layer supplies the explicit
pre-shift/shifted/backtranslation display form. The compatibility check preserves
physical and mathematical readout data. The normal form requires all new
active material to be shifted-normal or explicitly historical/pre-shift. Therefore the
release-final state is conservative over the prepared migration chain.
\end{proof}

\subsection{Release-release-final state non-claims}
\label{subsec:v400-release-final-non-claims}

\begin{remark}[\QSCver{shifted.0.0} release-final non-claims]
\label{rem:v400-release-final-non-claims}
The \QSCver{shifted.0.13} historical source-hygiene layer does not close OP~1/4 or OP~5(B), does not
activate primitive mixed tensors, does not claim Turing-completeness, and does
not replace standard QFT/QED/Lifshitz or general-relativistic calculations. It
is a global QAM/router/readout language migration and normal-form release
release-final state.
\end{remark}



\clearpage
